\subsection{Kafka Ingestion and Replay Design}

\subsubsection{Ingestion layer}
Events are published to Apache Kafka~\cite{kreps2011kafka}, running as a single
broker in Docker. There is one topic per source (\texttt{firms.events},
\texttt{epa.pm25.events} and \texttt{noaa.weather.events}), each with three
partitions and a replication factor of one. Each message is keyed by county name,
so all events for a county are kept in order within one partition. The message
value is a JSON object with a common envelope (source, county, event hour, event
identifier, replay scenario label, delay and scheduled arrival time) and a
source-specific payload. We checked that every event in each source has the same
set of fields.

\subsubsection{Replay input}
The input contains 18{,}000 events: 6{,}000 per source, covering five counties
at one-hour resolution over 50 days (1{,}200 hours) of the 2021 Dixie Fire period.
All event identifiers are unique. The replay producer reads these events and
sends each to the topic that matches its source.

\subsubsection{Replay speeds}
Replay speed is the number of simulated event-hours replayed per real second:
1 (slow), 10 (medium) and 100 (fast). Because the input spans 1{,}200 simulated
hours, a full replay takes about 1{,}200, 120 and 12 seconds at the three
settings. In a live test, the fast setting completed in 12 seconds.

\subsubsection{Delay conditions}
Late data is modelled with three delay levels. With no delay, every event arrives
on schedule. In the moderate condition, 1{,}846 of the 18{,}000 events (10.26\%)
are delayed by 1 to 3 hours. In the severe condition, 4{,}549 events (25.27\%) are
delayed by 4 to 12 hours. Each delay is recorded in the message as a delay in
hours, and the scheduled arrival time is the event hour plus that delay. The producer sends messages in order of scheduled arrival time and paces the replay by that time, so a delayed event is sent later than undelayed events from the same hour and can arrive after events with later event hours.

Events are selected for delay by a SHA-256 rule applied to the event identifier, with target fractions of 10\% (moderate) and 25\% (severe); the realised fractions are the 10.26\% and 25.27\% above. Delay selection is deterministic, depends only on the event identifier and the delay level, and is independent of both replay speed and the
scenario label: in a dry run over all 18{,}000 events, every speed delayed the
same events by the same amounts. Combining three speeds with three delay levels
gives nine experimental conditions, each intended to be repeated three times.

\subsubsection{Run isolation}
Because Kafka topics retain messages, repeating a run without clearing them would
mix events from different runs. Our experiment runner therefore deletes and
recreates the three topics before every run, checks that all topic offsets are
zero, replays the input, and records the duration, the producer exit status and
the number of messages in each topic.

\subsubsection{Verification}
A live replay of the full input sent all 18{,}000 events with no invalid records.
Each topic then held exactly 6{,}000 messages, and consumer samples from each
topic showed messages whose source field matched the topic.

Delay selection was also checked by running the moderate condition twice at each of the three speeds in dry-run mode. All six outputs of 18{,}000 messages had the same SHA-256 hash, so the messages, delays, arrival times and order were identical across runs and speeds. Live replays of the moderate and severe conditions delayed 1{,}846 and 4{,}549 events, matching the dry run, and the medium-speed replay took 120 seconds, as the speed setting predicts. A live slow-speed replay (one simulated event-hour per real second) took 1{,}200\,s and delivered all 18{,}000 messages with a clean producer exit.

\subsubsection{Technical Difficulty and Lesson Learned}
The first live slow-speed replay, which should take about 20 minutes, had not finished after two hours. The runner was still alive, but no result line had been written to the summary file and the producer log held only start-up warnings. The most likely cause was the laptop going to sleep, which pauses the timed waits that the replay depends on, so the measured duration would have been meaningless. We stopped the run, discarded its partial log, and repeated it with macOS sleep prevention enabled for the whole run (\texttt{caffeinate -dims}). The repeat finished in 1{,}200\,s with a clean producer exit and 6{,}000 messages in each topic. The lesson is that long paced replays must run on a machine that cannot sleep, and that every run must be judged by recorded evidence rather than by the process still running. The experiment runner therefore records the duration, the producer exit code and the per-topic message counts for every run, which makes a stalled or interrupted run easy to detect and exclude.

\subsubsection{Limitations}
The setup uses one broker with no replication, so it demonstrates the data flow
but not fault tolerance. Replay speeds are simulated and are not a measurement of
the original data rates.
